\chapter{Background and Related Work}
\label{chap:background}

This chapter introduces the ideas needed to understand the system without describing its implementation. It first explains computational reproducibility and the special properties of Jupyter notebooks. It then reviews earlier notebook studies, the three source platforms, metadata and provenance, environment reconstruction, classification, and semantic publication. The final section identifies the gap addressed by this thesis.

\section{Computational Reproducibility}

Computational results are reproducible when another person can use the same input data, code, computational steps, and analysis conditions and obtain consistent results. This thesis follows the terminology of the National Academies, which separates \emph{reproducibility} from \emph{replicability}: replication answers the same research question with newly collected data, whereas reproduction repeats the original computation \cite{nasem2019}. The distinction matters because this work evaluates digital research objects. It does not repeat the scientific study or judge whether its conclusion is correct.

Reproduction is not a single yes-or-no event. A rerun can stop while the software environment is being created, fail during execution, finish with changed outputs, or finish with outputs that match the stored notebook. Each outcome provides different evidence. A missing package points to incomplete dependency information. A network error may show that an external service is no longer available. A changed numerical result may arise from a newer library, random behaviour, or updated data. Therefore, a useful pipeline records the stage, error, duration, and comparison result instead of reducing every attempt to ``failed.''

Three groups of information influence computational reproducibility. The first is the \emph{computational object}: source code, notebook cells, input files, and stored outputs. The second is the \emph{environment}: programming-language version, libraries, operating-system assumptions, and external tools. The third is the \emph{provenance}: where the resource came from, when it was retrieved, what process acted on it, and which result that process produced. Access to code alone covers only part of this evidence.

Large notebook studies confirm the practical importance of these distinctions. Pimentel et al.~\cite{pimentel2019} examined quality and reproducibility across a large notebook corpus and connected successful execution with practices such as ordered cells and dependency information. Samuel and Mietchen~\cite{samuel2024computational} analysed notebooks linked to biomedical publications and recorded environment construction, execution exceptions, and output comparison separately. Their work shows that a pipeline must preserve intermediate outcomes if the results are to support later analysis.

Two operational definitions are used throughout this thesis. \emph{Successful execution} means that every selected code cell completes in the reconstructed environment without raising an uncaught exception. \emph{Output reproducibility} is stricter and is assessed per cell: the outputs of the freshly executed notebook are compared against the outputs stored in the published notebook, and each code cell is recorded as producing an identical result, a different result, or a result recognised as nondeterministic (for example a timestamp, a memory address, or a value derived from an unseeded random source). This cell-level comparison is the measurement introduced by Samuel and Mietchen~\cite{samuel2024computational}, and it is retained unchanged in this work; the per-notebook reproducibility score reported in Chapter~\ref{chap:evaluation} is the proportion of comparable code cells whose outputs matched.

Keeping the granularity at the cell level matters. A notebook in which one of forty cells diverges is a materially different outcome from a notebook in which none of the forty matched, and a binary reproduced/not-reproduced verdict cannot distinguish the two. These definitions make the evaluation measurable while acknowledging that identical output is not proof of scientific validity: they answer the narrower question of whether the published computation can still be rerun, and how much of its recorded result survives that rerun.

\clearpage
\section{Jupyter Notebooks and Reproducibility}

A Jupyter notebook is both a document and an executable program. It can place source code beside natural-language explanations, equations, tables, figures, and interactive results. This combination makes notebooks useful for exploring data and explaining an analysis to readers \cite{kluyver2016}. It also creates a challenge: the visible document does not necessarily contain everything needed to repeat the computation.

The standard \texttt{.ipynb} file is a JSON document containing an ordered list of cells and notebook-level metadata \cite{nbformat}. Code cells may store source code, an execution count, and several output types. Markdown cells store explanatory text, and raw cells can preserve content for later conversion. Kernel and language information may appear in metadata, but most metadata fields are optional. A valid notebook can therefore omit the exact Python or package versions used by its author.

Table~\ref{tab:notebook-anatomy} lists the notebook information and its value for reproduction.

\begin{table}[ht]
\centering
\caption{Notebook information and its value for reproduction.}
\label{tab:notebook-anatomy}
\begin{tabularx}{\textwidth}{@{}lXX@{}}
\toprule
Element & What it contains & Why it matters\tabularnewline
\midrule
Code cell & Program statements and imports & Reveals operations and undeclared packages\tabularnewline
Markdown cell & Explanations, headings, and instructions & May describe purpose, data, and execution steps\tabularnewline
Stored output & Text, images, tables, or errors & Provides a reference for comparison\tabularnewline
Execution count & The recorded order of code execution & Can reveal out-of-order interactive work\tabularnewline
Notebook metadata & Kernel and language hints & Helps select a compatible environment\tabularnewline
Repository files & Requirements, setup files, and data & Supplies context missing from the notebook\tabularnewline
\bottomrule
\end{tabularx}
\end{table}

Interactive execution is the central difficulty. A user may run cell 10, return to cell 3, redefine a variable, and then save the notebook. The displayed outputs can depend on hidden kernel state that is not reproduced by a clean top-to-bottom run. Rule et al.~\cite{rule2018} explain that notebooks are often used first for trying ideas and exploring data. During this work, users may run cells in different orders or test several versions of an analysis. The saved notebook may therefore not clearly show how the final result was produced. Kery et al.~\cite{kery2018} found that data scientists often organize and explain the notebook only after the exploratory work is complete.

Reproducible execution therefore requires more than opening the file. The notebook should be run in a new kernel, in a known order, with explicit dependencies and available data. Stored outputs should be preserved long enough to compare them with new outputs. The surrounding repository is equally important because it may contain local modules, requirements files, setup scripts, or documentation. This thesis treats the notebook as the central executable object while using its repository or archived record as the context needed to interpret and rerun it.

\clearpage
\section{Existing Notebook Reproducibility Analysis}

Earlier studies examine notebook reproducibility in different ways. Some check the quality of many notebooks. Others try to execute notebooks automatically or organize the results in a knowledge graph. These studies provide the starting point for this thesis.

Pimentel et al.~\cite{pimentel2019} studied more than one million notebooks from GitHub and reran a selected group under controlled conditions. They found that reproducibility is related to features such as a clear execution order, version control, and documented dependencies. Samuel and Mietchen~\cite{samuel2024computational} later studied notebooks connected to biomedical publications. Their pipeline found the GitHub repositories, created Python environments, executed the notebooks, recorded errors, and compared the new outputs with the saved outputs. The results were stored as structured database records, not only as downloaded files.

FAIR Jupyter then converted this database information into a knowledge graph \cite{samuel2024fairjupyter}. This allows users to ask connected questions about publications, repositories, notebooks, executions, and results. However, both the source pipeline and the first graph version mainly describe GitHub resources. Other research gives practical advice to notebook authors. Rule et al.~\cite{rule2019tenrules}, for example, recommend a clear cell order, documented dependencies, and visible data-processing steps.

Table~\ref{tab:related-work} summarises the relationship between key prior work and this thesis.

\begin{table}[ht]
\centering
\caption{Relationship between key prior work and this thesis.}
\label{tab:related-work}
\small
\begin{tabularx}{\textwidth}{@{}p{2.6cm}XXX@{}}
\toprule
Work & Main focus & Useful foundation & Remaining limitation\tabularnewline
\midrule
Pimentel et al.~\cite{pimentel2019} & Large-scale notebook quality and reruns & Quality indicators and execution evidence & GitHub corpus; no cross-platform connector model\tabularnewline
Samuel and Mietchen~\cite{samuel2024computational} & Biomedical notebook reproduction & Baseline database, environment, errors, comparisons & Resource acquisition centred on GitHub\tabularnewline
FAIR Jupyter~\cite{samuel2024fairjupyter} & Semantic exploration of results & Ontology, mappings, provenance-oriented queries & GitHub repository representation\tabularnewline
This thesis & Cross-platform reproducibility & Reuses execution and semantic foundations & ---\tabularnewline
\bottomrule
\end{tabularx}
\end{table}

This thesis builds on these works~\cite{pimentel2019,samuel2024computational,samuel2024fairjupyter} instead of replacing them. The existing environment, execution, and comparison stages already record useful information and are therefore reused. The new work starts when a resource enters the pipeline. The system must now identify, check, describe, and download resources from three platforms. It must also record the platform in the database, classification results, and knowledge graph. This makes it possible to compare metadata and reproducibility across platforms, which a GitHub-only dataset cannot do.

\clearpage
\section{GitHub, Codeberg, and Zenodo}

GitHub, Codeberg, and Zenodo can all publish notebooks, but they store different kinds of resources. GitHub and Codeberg are Git hosting platforms. Their main resource is a version-controlled repository. Zenodo is a research repository. Its main resource is a published record with metadata and downloadable files. The pipeline must keep this difference because one download method cannot correctly handle all three platforms.

GitHub provides two forms of access. Git is used to clone the files, while the REST API provides information such as the name, owner, description, topics, licence, and dates \cite{githubapi}. Codeberg also supports Git, but it uses the Forgejo software and API \cite{codebergdocs,forgejoapi}. Therefore, GitHub and Codeberg can be cloned in a similar way, but their API addresses and JSON fields are not always the same~\cite{codebergapi}. Their notebook links also differ: GitHub normally uses \texttt{/blob/}, while Codeberg uses \texttt{/src/}.

Zenodo gives each published record a numeric identifier and, on publication, a DOI~\cite{zenodoapi}. Its Records API returns research metadata and a list of files that can be downloaded \cite{zenodoapi}. A record may contain notebooks, source-code archives, data, documentation, or several of these. It normally has no Git remote. Therefore, the Zenodo part of the pipeline must check the API record, select a useful file, download it, extract it, and then search the extracted folder for notebooks.

Table~\ref{tab:platform-differences} shows the platform properties that affect reproducibility processing.

\begin{table}[ht]
\centering
\caption{Platform properties that affect reproducibility processing.}
\label{tab:platform-differences}
\small
\begin{tabularx}{\textwidth}{@{}lXXX@{}}
\toprule
Aspect & GitHub & Codeberg & Zenodo\tabularnewline
\midrule
Primary object & Git repository & Git repository & Archived research record\tabularnewline
Stable local acquisition & Clone or pull & Clone or pull & Download and extract\tabularnewline
Validation signal & Git remote & Git remote & Records API\tabularnewline
Metadata source & GitHub REST API & Forgejo REST API & Zenodo Records API\tabularnewline
Identifier & owner/repository & owner/repository & numeric record ID and DOI\tabularnewline
Version context & commits and branches & commits and branches & record versions and deposited files\tabularnewline
\bottomrule
\end{tabularx}
\end{table}

These differences motivate separating platform-specific acquisition from platform-independent notebook processing. Only the early steps depend on the platform: confirming that a resource exists, reading its description, and obtaining its files. Once those files are available locally, the notebooks they contain can be discovered, executed, and compared in the same way regardless of where they came from. A study spanning several platforms therefore does not require a separate analysis workflow for each of them. The connector architecture adopted in this thesis is presented in Section~\ref{sec:connector-design}.

\section{Metadata Normalization and Provenance}

Metadata and provenance answer two different questions. Metadata answers, ``What is this resource?'' It includes information such as the title, authors, description, licence, keywords, DOI, and dates. Provenance answers, ``Where did the resource come from, and what did the pipeline do with it?'' For example, metadata can say that a Zenodo record has a certain title and two authors. Provenance can show that the record came from Zenodo, was downloaded by the pipeline, and was processed during a specific run.

GitHub, Codeberg, and Zenodo do not return this information in the same form. GitHub may provide one owner account, while Zenodo may provide a list of creators. One platform may return a licence as a short name, another may return a larger data object, and sometimes no licence is available. These differences make direct comparison difficult.

Metadata normalization is the general technique for resolving such differences. It maps the differing responses of several sources onto one shared set of descriptive elements, so that values with the same meaning can be compared regardless of where they came from. Which elements a study selects depends on the questions it asks, but established metadata standards provide a common starting point: the DataCite schema, for example, defines a small mandatory core of title, creator, publisher, publication year, and identifier for describing research objects \cite{datacite45}.

A normalization scheme has to decide what to do when a source provides no value for an element. The decision matters more than it first appears. A scheme that fills the gap, by substituting a default or by deriving a value from elsewhere, produces a uniformly complete description in which a genuinely absent value can no longer be distinguished from a supplied one. A scheme that records the absence keeps that distinction, and thereby makes metadata completeness itself measurable. For a study comparing how well different platforms describe what they host, the second choice is the only workable one.

Dates require the same care. A date reported by a source describes the resource, recording when its author created or last changed it. A date recorded by an analysis describes the analysis, recording when the resource was retrieved or processed. The two belong to different histories, and conflating them would make the record of the study indistinguishable from the record of the object being studied.

The knowledge graph records this history using PROV-O, a standard vocabulary for provenance \cite{prov2013}. It has three main ideas. An \texttt{Entity} is a resource, such as a repository, notebook, or result. An \texttt{Activity} is an action, such as a repository run or notebook execution. An \texttt{Agent} is responsible for a resource or action, such as GitHub, Codeberg, or Zenodo. For example, Zenodo is the agent that provides a record, the record is an entity used by the pipeline run, and the run is an activity that produces an execution result.

Together, normalization and provenance make results from heterogeneous sources both comparable and checkable. Normalization makes equivalent values comparable; provenance keeps visible the path by which each value was obtained, so that a surprising result can be traced back to its origin rather than merely accepted or discarded. The concrete normalization model adopted in this thesis is presented in Section~\ref{sec:metadata-model-design}, and the provenance model in Section~\ref{sec:kg-design}.

\clearpage
\section{Environment Reconstruction and Notebook Execution}

A notebook needs more than the code stored in its cells. It may depend on a specific Python version, external packages, local modules, system programs, input data, environment variables, or online services. Environment reconstruction means using the available files and notebook content to create an environment in which the notebook can run.

Files such as \texttt{requirements.txt}, \texttt{setup.py}, and Python-version files provide the clearest dependency information. When these files are missing or incomplete, imports in notebook cells provide additional clues. However, an imported module can have a different name from the package installed with \texttt{pip}, and an import normally gives no version. Automatically generated requirements are therefore only the pipeline's best attempt to rebuild the environment. They are not an exact copy of the author's original computer.

Selecting an appropriate interpreter version is equally important. Code written for one Python version can fail under another, and many packages are only compatible with a bounded version range. Version-management tools such as \texttt{pyenv} allow several interpreter releases to coexist on one machine and to be selected per project \cite{pyenv}, and installing each project's packages into a separate virtual environment prevents the dependencies of one resource from perturbing another. Wang et al.~\cite{wang2020restoring} show that automatically restoring a notebook's original execution environment is itself a hard problem, and that even careful reconstruction frequently fails to recover the author's setting exactly.

Jupyter \texttt{nbconvert} executes notebooks automatically and saves the new outputs \cite{nbconvert}. Each run starts with a new kernel and executes the cells from top to bottom. This removes variables left over from an earlier interactive session. It can also reveal hidden assumptions. For example, a later cell may need a variable that the author created by running cells in a different order. Chattopadhyay et al.~\cite{chattopadhyay2020} found that environment setup, dependency management, and turning exploratory notebooks into reliable workflows are common problems.

Execution success and output reproducibility are different results. A notebook can finish without errors but still produce different numbers, text, tables, or figures. This can happen because of random behaviour, current dates, changed online data, or new package versions. The comparison stage records these changes instead of treating every difference as a simple error. Matching outputs also do not prove that the scientific method is correct. They only show that the tested run produced the same saved result.

For this reason, a reproducibility study must treat re-execution as a staged process and observe each stage separately rather than collapsing an attempt into a single success-or-failure verdict. Dependency resolution, interpreter selection, package installation, execution, and output comparison can each fail for different reasons and carry different evidential meaning. Studies that preserve these intermediate outcomes \cite{pimentel2019,samuel2024computational} can attribute a failure to a specific cause, whereas studies that record only a final verdict cannot.

\clearpage
\section{Notebook Classification Approaches}

Reproducibility results are easier to understand when notebooks are grouped by their main purpose. A tutorial, a data-cleaning notebook, and a machine-learning experiment can need different packages and may fail for different reasons. Classification helps separate these differences from differences caused by the hosting platform.

A notebook does not normally contain one field that clearly states its purpose. Useful clues appear in the filename, headings, Markdown explanations, imports, function calls, outputs, and order of cells. Earlier studies show that notebooks often mix exploration with later explanation \cite{rule2018,kery2018}. The classifier must therefore examine both written text and code. It must also allow a notebook to have more than one important purpose.

Table~\ref{tab:classification-approaches} gives the classification strategies considered in this thesis.

\begin{table}[ht]
\centering
\caption{Classification strategies considered in this thesis.}
\label{tab:classification-approaches}
\small
\begin{tabularx}{\textwidth}{@{}p{2.7cm}XXX@{}}
\toprule
Approach & Strength & Limitation & Research relevance\tabularnewline
\midrule
Weighted rules & Transparent and repeatable; works without a model & Keywords and imports can be ambiguous & Transparent baseline\tabularnewline
AI model & Can combine code and narrative context & Output may vary and requires validation & Context-sensitive alternative\tabularnewline
Human review & Can resolve context and mixed intent & Slow and not suitable for every item & Reference for ambiguous cases and validation\tabularnewline
Hybrid decision & Exposes agreement and uncertainty & Requires a clear conflict policy & Supports analysis of agreement and uncertainty\tabularnewline
\bottomrule
\end{tabularx}
\end{table}

The categories used for such a characterisation are not arbitrary. Corpus studies of notebook content converge on a recurring set of activities. Rule et al.~\cite{rule2018} distinguish exploratory analysis from explanatory and instructional material. Kery et al.~\cite{kery2018} describe an iterative cycle in which data is first prepared, then analysed, and only later documented. Pimentel et al.~\cite{pimentel2019} separate notebooks by their dominant library usage, which cleanly distinguishes data-manipulation work from modelling work and from plotting work. Psallidas et al.~\cite{psallidas2022datascience}, analysing millions of GitHub notebooks, report the same broad activity groups --- data acquisition and preparation, exploration and analysis, visualization, and model training --- as the dominant uses of the format. Quaranta et al.~\cite{quaranta2021kgtorrent} observe an equivalent distribution in a large Kaggle corpus, and Koenzen et al.~\cite{koenzen2020duplication} additionally identify tutorial and template notebooks as a distinct population that is copied and adapted rather than authored from scratch.

The taxonomy adopted in this thesis is derived from these sources and is defined in Section~\ref{sec:classification-taxonomy}. Two observations from the literature constrain any such scheme. First, a notebook frequently serves several purposes at once, because preparation, analysis, and visualization commonly occur within one document \cite{rule2018,kery2018}; a usable scheme therefore requires an explicit convention for multi-purpose notebooks rather than forcing a single label. Second, the evidence available in a notebook is sometimes too sparse to support any confident assignment, so a scheme also requires a way to express that the evidence is insufficient instead of guessing.

These limitations motivate combining complementary automated approaches rather than relying on a single one, and reserving human review for the cases where the automated evidence is weak or contradictory. Research on human involvement in machine learning supports reviewing unclear cases instead of accepting every model answer as correct \cite{mosqueira2023hitl}. Whichever approaches are combined, their reliability cannot be established from their agreement with one another: two methods can agree and both be wrong, so a comparison against labels assigned by a person remains necessary. The classification architecture adopted in this thesis is presented in Section~\ref{sec:classification-design}.

\clearpage
\section{FAIR Data and Knowledge Graphs}

The FAIR principles state that digital research objects should be findable, accessible, interoperable, and reusable \cite{wilkinson2016fair}. They apply to data, code, workflows, and other research materials. FAIR information should be understandable not only to people but also to software. For notebook reproducibility, providing a download link is not enough. The notebook should be connected to its platform, metadata, processing history, execution, and result in a machine-readable form.

RDF stores information as simple statements called triples. Each triple has a subject, a relation, and an object \cite{rdf2014}. For example, ``repository 42 -- is hosted on -- Codeberg'' is one triple. ``run 17 -- processed -- repository 42'' is another. Because both statements use the same repository identifier, they become connected in the graph. SPARQL is a query language that follows these connections \cite{sparql2013}. It can move from a platform to its repositories, notebooks, executions, and results without manually joining several CSV files.

An ontology defines the types of things and the relations allowed in the graph. FAIR Jupyter reuses established vocabularies, among them PROV-O for processing history and DOAP for software projects, and adds terms specific to notebook reproducibility \cite{samuel2024fairjupyter}. Its published graph describes publications, GitHub repositories, notebooks, executions, and reproducibility results drawn from the biomedical study. Its vocabulary is correspondingly scoped: it models the Git repository as the unit of publication, and has no terms for a hosting platform as such, for an archived research record, or for the processing activity that produced a given result. Any study spanning more than one kind of publication venue therefore requires an extension of this model.

Graphs of this kind are normally not written by hand. Declarative mapping languages describe how records in an existing source --- a relational table or a CSV export --- become RDF statements. RML generalises the relational R2RML standard to heterogeneous sources \cite{rmlspec}, and YARRRML provides a human-readable surface syntax that compiles into RML \cite{heyvaert2018yarrrml}. A materialisation engine such as Morph-KGC \cite{arenas2024morph} then reads the mappings and emits the triples. The significance of this arrangement for reproducibility research is that the ontology, the mapping rules, and the resulting graph remain three separate, independently inspectable artefacts: a reviewer can examine the model and the transformation rules, and the graph can be regenerated whenever the underlying data changes.

A materialised graph also has to be validated. Validation normally combines structural checks --- that the ontology and mappings parse, and that source records yield the expected entities and relations --- with competency questions, that is, SPARQL queries formulating the questions the graph was built to answer. A graph that parses but cannot answer its competency questions represents a large file of triples rather than a usable model.

\clearpage
\section{Research Gap}
\label{sec:research-gap}

The work reviewed in this chapter establishes that notebook re-execution can be automated at scale \cite{pimentel2019,samuel2024computational}, that its outcomes can be recorded at cell-level granularity \cite{samuel2024computational}, and that the resulting evidence can be published semantically and queried \cite{samuel2024fairjupyter}. What it does not establish is whether any of these findings hold beyond the single platform on which they were obtained. Four gaps follow, each corresponding to one research question stated in Section~\ref{sec:research-questions}.

\paragraph{Gap 1: platform-independence has been assumed rather than demonstrated.} The execution, comparison, and error-classification stages of existing pipelines are plausibly independent of where a notebook came from, but this has never been tested, because every published corpus originates from GitHub. The boundary at which platform-specific handling genuinely ends is therefore unknown, as is whether the invariant holds for an archival record --- an object that is not a Git repository and that may package its notebooks inside an archive without any accompanying environment specification. Establishing where that boundary lies is the subject of RQ1.

\paragraph{Gap 2: the representational cost of metadata heterogeneity is unmeasured.} GitHub, Forgejo, and Zenodo describe a resource through incompatible vocabularies with different mandatory fields \cite{githubapi,forgejoapi,zenodoapi}. Integrating them requires a normalized representation, but normalization is lossy, and the direction of the loss matters: a scheme that silently fills absent values destroys exactly the signal that distinguishes a well-described archival deposit from a bare repository. No existing study reports how much descriptive information survives such an integration, nor how completeness varies by venue. This is the subject of RQ2. The FAIR Jupyter ontology \cite{samuel2024fairjupyter} provides a suitable semantic foundation but models GitHub repositories only; it has no vocabulary for a hosting platform, an archived research record, a pipeline run, or a normalized description.

\paragraph{Gap 3: notebook purpose is unmeasured, and the available automated methods have not been compared.} The corpus studies discussed in Section~2.7 agree on the broad activity types that notebooks serve, but none attaches a purpose label to the individual notebooks whose reproducibility is being measured. Two families of automated labelling are available: transparent rules over imports and text, and language models that read the notebook as a document. They rest on different evidence and can be expected to fail in different circumstances, yet how far they converge on the same description, and whether their weaknesses are independent enough to make combining them worthwhile, has not been reported. This is the subject of RQ3.

\paragraph{Gap 4: cross-platform differences are unknown and currently unmeasurable.} Because Gaps 1--3 are open, the empirical question cannot presently be asked: it is not known whether metadata completeness, execution outcomes, failure causes, or output agreement differ between a commercial forge, a non-profit forge, and an archival repository. A comparison drawn without controlling for notebook purpose would additionally risk attributing a composition effect to the platform. Closing this gap requires the instrument that Gaps 1--3 describe, and is the subject of RQ4.

These four gaps are addressed by one system rather than four separate artefacts, because they are interdependent: the empirical comparison requires the integrated metadata, which requires the platform-independent workflow, and its interpretation requires validated purpose labels. Chapter~\ref{chap:requirements} turns these gaps into verifiable requirements, Chapter~\ref{chap:design} presents the design that satisfies them, Chapter~\ref{chap:implementation} describes its realisation, and Chapter~\ref{chap:evaluation} reports the resulting evidence.

The scope of what such an instrument can settle is bounded, and the bound should be stated here rather than discovered later. A sample drawn from three platforms cannot represent the full population of notebooks on any of them. A reconstructed environment is an approximation of the author's original setting, not a recovery of it \cite{wang2020restoring}. An automated purpose label is a prediction, not a ground truth. What the system can provide is a consistent, traceable, and repeatable procedure that applies identical treatment to resources of heterogeneous origin, and that keeps the origin of every measurement visible in the result.
